\NSPage{097}%
On the reserved mean patch, the background retains the power law \eqref{eq:4.30} by
Proposition~\ref{prop:5.5}. We will use this form to reduce the equations to independent power
moments of supported bumps. We first construct the fixed linear map for arbitrary targets
(Lemma~\ref{lem:8.7}), then compute what remains when its targets are the current defects
(Lemma~\ref{lem:8.8}).

In \NSi{NS-F-P097-001}{x=r/\sqrt q\in I_m}, the base after velocity normalization at
\NSi{NS-F-P097-002}{q} is
\NSFormula{NS-F-P097-003}{display}{8.24}
\begin{equation}
G_q=0,\qquad V_q=a(\eta)x^{-1-2\lambda},\qquad |a(\eta)|\ge a_0>0.
\tag{8.24}\label{eq:8.24}
\end{equation}
Here \NSi{NS-F-P097-004}{a(\eta)=2^{1/2+\lambda}c_{\mathrm{patch}}(1+\eta^2)^{-1}}. The exact
summed-base statement is \eqref{eq:5.44}. The tangential coefficients beyond the leading profile
are supported elsewhere by \eqref{eq:5.18}. Multiplication of their azimuthal potentials by a cutoff
in \NSi{NS-F-P097-005}{q} does not create an axial component on this patch, because
\NSi{NS-F-P097-006}{q} is independent of radius. Thus \eqref{eq:8.24} holds for the summed base
used in the balances. The fixed parameter \NSi{NS-F-P097-007}{\lambda} is positive.

\begin{lemma}\label{lem:8.7}
Suppose that the fixed slow base satisfies \eqref{eq:8.24} on the reserved interior interval
\NSi{NS-F-P097-008}{I_m}, with \NSi{NS-F-P097-009}{\lambda>0} and
\NSi{NS-F-P097-010}{|a(\eta)|\ge a_0>0}. There are three fixed azimuthal bump profiles and two
fixed axial bump profiles supported in this interval with the following property. At each
\NSi{NS-F-P097-011}{(Z,T)}, for arbitrary scalar targets
\NSi{NS-F-P097-012}{(\mathcal P,\mathcal J_\theta,\mathcal J_z)}, their physical rescalings have a
unique linear combination \NSi{NS-F-P097-013}{(\Delta v,\gamma_d)} satisfying
\NSFormula{NS-F-P097-014}{display}{8.25}
\begin{equation}
\begin{gathered}
\int R^2\Delta v\,dR=0,\qquad \int R\gamma_d\,dR=0,\\
\int (2V/R)\Delta v\,dR=-\mathcal P,\\
\int R^2(G\Delta v+V\gamma_d)\,dR=-\mathcal J_\theta,\\
\int (2RG\gamma_d-RV\Delta v)\,dR=-\mathcal J_z.
\end{gathered}
\tag{8.25}\label{eq:8.25}
\end{equation}
The target triple is normalized row by row as in \eqref{eq:8.17}. The resulting increments are
independent of the angular and auxiliary variables and depend linearly on the targets. For every
\NSi{NS-F-P097-015}{\alpha}, targets in \NSi{NS-F-P097-016}{\mathsf S_\alpha} give increments in
\NSi{NS-F-P097-017}{\mathsf M_\alpha}, for every fixed coefficient derivative, with constants permitted to
depend on the fixed \NSi{NS-F-P097-018}{\lambda}. The azimuthal-potential realization
\eqref{eq:8.14} has \NSi{NS-F-P097-019}{\Delta\gamma=\gamma_d} exactly and
\NSi{NS-F-P097-020}{\Delta\beta\in \mathsf M_{\alpha+1}}. The two integral constraints are preserved
exactly. When the targets are the current defects defined by Equations~\eqref{eq:8.12} and
\eqref{eq:8.15}, the last three equations cancel the displayed contributions involving the fixed
base.
\end{lemma}

\begin{proof}
\textbf{Step 1: choose the five profiles.} Normalize every physical row and its target using
\NSi{NS-F-P097-021}{q}. On the patch, \eqref{eq:8.24} splits the rows into an angular block whose
powers of \NSi{NS-F-P097-022}{x} are \NSi{NS-F-P097-023}{2,-2-2\lambda,-2\lambda}, and an axial
block whose powers are \NSi{NS-F-P097-024}{1,1-2\lambda}. Each list contains distinct powers for
the fixed positive \NSi{NS-F-P097-025}{\lambda}.

To see invertibility explicitly, choose a nonnegative nonzero bump
\NSi{NS-F-P097-026}{\eta_0} supported near some \NSi{NS-F-P097-027}{x_0>0} inside the patch.
Choose a small fixed \NSi{NS-F-P097-028}{d>0} and a sufficiently narrow support so that its
geometrically scaled copies
\NSi{NS-F-P097-029}{\eta_j(x)=a_j^{-1}\eta_0(x/a_j),\ a_j=e^{jd},\ j=0,1,2}, lie on disjoint
subintervals of the patch. Their power moments are
\NSFormula{NS-F-P097-030}{display}{none}
\[
\int x^p\eta_j(x)\,dx=\mu_p e^{jdp},\qquad
\mu_p=\int x^p\eta_0(x)\,dx>0.
\]
For the three angular powers the matrix, after dividing its rows by
\NSi{NS-F-P097-031}{\mu_p}, is the ordinary Vandermonde matrix in the distinct numbers
\NSi{NS-F-P097-032}{e^{dp}}. Its determinant is nonzero. Two copies give the same argument for
the axial block. The remaining factors are nonzero multiples of
\NSi{NS-F-P097-033}{a(\eta)}.

\NSPage{098}%
\textbf{Step 2: define the coefficients by matrix inversion.} Use
\NSi{NS-F-P098-001}{\eta_0,\eta_1,\eta_2} as the azimuthal profiles and
\NSi{NS-F-P098-002}{\eta_0,\eta_1} as the axial profiles. Let
\NSi{NS-F-P098-003}{\mathcal P_q,\mathcal J_{\theta,q},\mathcal J_{z,q}} denote the physical
targets multiplied by the powers in \eqref{eq:8.17}, with \NSi{NS-F-P098-004}{q} in place of
\NSi{NS-F-P098-005}{Q}. Define the matrices
\NSFormula{NS-F-P098-006}{display}{none}
\[
\begin{gathered}
A_\theta=
\begin{pmatrix}
\mu_2 & \mu_2e^{2d} & \mu_2e^{4d}\\
2a\mu_{-2-2\lambda} & 2a\mu_{-2-2\lambda}e^{d(-2-2\lambda)}
  & 2a\mu_{-2-2\lambda}e^{2d(-2-2\lambda)}\\
-a\mu_{-2\lambda} & -a\mu_{-2\lambda}e^{-2d\lambda}
  & -a\mu_{-2\lambda}e^{-4d\lambda}
\end{pmatrix},
\\[4pt]
A_z=
\begin{pmatrix}
\mu_1 & \mu_1e^d\\
a\mu_{1-2\lambda} & a\mu_{1-2\lambda}e^{d(1-2\lambda)}
\end{pmatrix},
\qquad a=a(\eta).
\end{gathered}
\]
Step 1 proves that both matrices are invertible. Define
\NSFormula{NS-F-P098-007}{display}{none}
\[
\begin{gathered}
(u_0,u_1,u_2)^T=A_\theta^{-1}(0,-\mathcal P_q,-\mathcal J_{z,q})^T,\\
(s_0,s_1)^T=A_z^{-1}(0,-\mathcal J_{\theta,q})^T,\\
\Delta v_{\mathrm{phys}}(r,z,t)=q^{-A}\sum_{j=0}^{2}u_j\eta_j(r/\sqrt q),\\
\gamma_{d,\mathrm{phys}}(r,z,t)=q^{-A}\sum_{j=0}^{1}s_j\eta_j(r/\sqrt q).
\end{gathered}
\]
The matrix \NSi{NS-F-P098-008}{A_\theta} enforces the zero angular-momentum constraint and the
targets \NSi{NS-F-P098-009}{-\mathcal P_q,-\mathcal J_{z,q}};
\NSi{NS-F-P098-010}{A_z} enforces the zero axial-flux constraint and the target
\NSi{NS-F-P098-011}{-\mathcal J_{\theta,q}}. Since \NSi{NS-F-P098-012}{G_q=0}, these two
blocks exhaust the five equations. Their inverse defines one physical correction, linear in the
target triple and smooth in the slow variables whenever the targets are smooth. The inverse may
deteriorate as \NSi{NS-F-P098-013}{\lambda\downarrow0}; no uniformity in that limit is required.

\textbf{Step 3: normalization, bounds, and axial realization.} Passing from
\NSi{NS-F-P098-014}{q}-normalized rows to a fixed \NSi{NS-F-P098-015}{Q} chart multiplies each
row and its target by the same powers in \eqref{eq:8.17} and by bounded smooth powers of
\NSi{NS-F-P098-016}{q/Q}. Those ratios have bounded derivatives of every fixed order. Matrix
inversion and differentiation therefore cost fixed constants, and the radial weight is bounded above
and below on the fixed interior support. This proves the \NSi{NS-F-P098-017}{\mathsf S_\alpha} to
\NSi{NS-F-P098-018}{\mathsf M_\alpha} bound. The desired axial field is slow and its weighted axial-flux
integral is zero; Lemma~\ref{lem:8.2} gives an identically vanishing cutoff remainder. Its azimuthal
potential and radial velocity component stay in the mean patch, including between the finitely many
bump supports.
\end{proof}

The linear map has now been fixed independently of the target values. Its use in the nonlinear
equations requires the defects of the updated field. The next statement records the exact difference
and the gain available under the bounds used in the iteration. The base estimates for this gain are
local to the support of the azimuthal potential, which includes the intervals between its bump
supports.

\begin{lemma}\label{lem:8.8}
Suppose that the slow base satisfies \eqref{eq:8.24}, with the fixed
\NSi{NS-F-P098-019}{\lambda>0} and lower bound
\NSi{NS-F-P098-020}{|a(\eta)|\ge a_0>0}. Let
\NSi{NS-F-P098-021}{(\beta,v,\gamma)} be a smooth divergence-free mean correction supported in
the active shell, and keep the base and wave field \NSi{NS-F-P098-022}{w} fixed. Form
\NSi{NS-F-P098-023}{g_r} and
\NSi{NS-F-P098-024}{(\mathcal P,\mathcal J_\theta,\mathcal J_z)} by Equations~\eqref{eq:8.3},
\eqref{eq:8.12} and \eqref{eq:8.15}. Apply Lemma~\ref{lem:8.7} with these three defects as its
targets, and realize the resulting \NSi{NS-F-P098-025}{\gamma_d} by \eqref{eq:8.14}. For the
updated mean velocity
\NSFormula{NS-F-P098-026}{display}{none}
\[
(\beta_{\mathrm{new}},v_{\mathrm{new}},\gamma_{\mathrm{new}})
=(\beta+\Delta\beta,v+\Delta v,\gamma+\Delta\gamma),
\]
\NSPage{099}%
compute \NSi{NS-F-P099-001}{g_{r,\mathrm{new}}} from \eqref{eq:8.3} and define
\NSi{NS-F-P099-002}{\mathcal R_g=g_{r,\mathrm{new}}-g_r-(2V/R)\Delta v}. Then
\NSFormula{NS-F-P099-003}{display}{8.26}
\begin{equation}
\begin{aligned}
\mathcal R_g={}&-\mathfrak t_*\Delta\beta
 -(D_r+1/R)\bigl(2b\Delta\beta+2\beta\Delta\beta+(\Delta\beta)^2\bigr)\\
&-D_z\bigl(b\Delta\gamma+G\Delta\beta+\beta\Delta\gamma
  +\gamma\Delta\beta+\Delta\beta\Delta\gamma\bigr)\\
&+(2v\Delta v+(\Delta v)^2)/R+\varepsilon(\Delta_0-R^{-2})\Delta\beta.
\end{aligned}
\tag{8.26}\label{eq:8.26}
\end{equation}
The defects recomputed from the updated field satisfy
\NSFormula{NS-F-P099-004}{display}{8.27}
\begin{equation}
\begin{aligned}
\mathcal P_{\mathrm{new}}&=\int\langle\mathcal R_g\rangle_Y\,dR,\\
(\mathcal J_\theta)_{\mathrm{new}}&=
 \int R^2\langle\gamma\Delta v+v\Delta\gamma+\Delta\gamma\Delta v\rangle_Y\,dR,\\
(\mathcal J_z)_{\mathrm{new}}&=
 \int R\langle2\gamma\Delta\gamma+(\Delta\gamma)^2\rangle_Y\,dR
 -\frac12\int R^2\langle\mathcal R_g\rangle_Y\,dR.
\end{aligned}
\tag{8.27}\label{eq:8.27}
\end{equation}
Suppose in addition that, on the support of the azimuthal potential for this correction, every fixed
coefficient derivative of \NSi{NS-F-P099-005}{b} is bounded by
\NSi{NS-F-P099-006}{C_I\varepsilon S_*^{B_I}}, and those of the slow
\NSi{NS-F-P099-007}{V,G} by \NSi{NS-F-P099-008}{C_I S_*^{B_I}}. If the current corrections
satisfy \NSi{NS-F-P099-009}{v,\gamma\in \mathsf M_{0.9}},
\NSi{NS-F-P099-010}{\beta\in \mathsf M_{1.9}}, and the three targets belong to
\NSi{NS-F-P099-011}{\mathsf S_\alpha} for \NSi{NS-F-P099-012}{\alpha\ge0.9}, then
\NSFormula{NS-F-P099-013}{display}{none}
\[
\mathcal R_g\in \mathsf M_{\alpha+0.9-2\kappa_s},\qquad
(\mathcal P_{\mathrm{new}},(\mathcal J_\theta)_{\mathrm{new}},
 (\mathcal J_z)_{\mathrm{new}})\in \mathsf S_{\alpha+0.9-2\kappa_s}.
\]
\end{lemma}

\begin{proof}
Subtract the two radial sources in \eqref{eq:8.3}. The covariance
\NSi{NS-F-P099-014}{W} is unchanged, and expanding each quadratic product gives
\eqref{eq:8.26}. In the azimuthal flux defect, the terms linear in the base are
\NSi{NS-F-P099-015}{\int R^2(G\Delta v+V\Delta\gamma)\,dR}. In the axial flux defect they are
\NSi{NS-F-P099-016}{\int(2RG\Delta\gamma-RV\Delta v)\,dR}, where the second term comes from
the radial source contribution
\NSi{NS-F-P099-017}{-\frac12\int R^2(2V/R)\Delta v\,dR}. Since
\NSi{NS-F-P099-018}{\Delta\gamma=\gamma_d}, the last three rows of \eqref{eq:8.25} cancel the old
defects and these linear terms. The remaining products give \eqref{eq:8.27}.

For the estimates, Lemma~\ref{lem:8.7} gives
\NSi{NS-F-P099-019}{\Delta v,\Delta\gamma\in \mathsf M_\alpha} and
\NSi{NS-F-P099-020}{\Delta\beta\in \mathsf M_{\alpha+1}}. All three increments are independent of the
auxiliary variable. In particular,
\NSi{NS-F-P099-021}{\mathfrak t_*\Delta\beta=-\varepsilon\partial_T\Delta\beta}. The terms in
\NSi{NS-F-P099-022}{\mathcal R_g} satisfy
\NSFormula{NS-F-P099-023}{display}{none}
\[
\begin{array}{c|c}
\text{term}&\text{class}\\ \hline
-\mathfrak t_*\Delta\beta&\mathsf M_{\alpha+2}\\
(D_r+1/R)(2b\Delta\beta)&\mathsf M_{\alpha+2-\kappa_s}\\
D_z(b\Delta\gamma+G\Delta\beta)&\mathsf M_{\alpha+2}\\
2v\Delta v/R&\mathsf M_{\alpha+0.9}\\
(\Delta v)^2/R&\mathsf M_{2\alpha}\\
\varepsilon(\Delta_0-R^{-2})\Delta\beta&\mathsf M_{\alpha+2-2\kappa_s}.
\end{array}
\]
The remaining transport products lie in \NSi{NS-F-P099-024}{\mathsf M_{\alpha+2-\kappa_s}}. Since
\NSi{NS-F-P099-025}{\alpha\ge0.9}, every exponent above is at least
\NSi{NS-F-P099-026}{\alpha+0.9-2\kappa_s}. Radial integration preserves these estimates after
normalizing each moment by the factors specified above. The radial weight is bounded below on the
interior support of the azimuthal potential, so the stated local base bounds suffice.
\end{proof}

\subsection{Recomputation and compatibility between charts.}\label{sec:8.7}
In this subsection we specify the order in which the residuals and defects in
Equations~\eqref{eq:8.3}, \eqref{eq:8.12} and \eqref{eq:8.15} are recomputed after a mean
correction, and verify that the constructions agree between charts using Equations~\eqref{eq:8.4}
and \eqref{eq:8.23} and Lemma~\ref{lem:6.2}. Each correction changes the quadratic fluxes, so the
next source must be computed from the updated divergence-free velocity. A cutoff remainder with zero
auxiliary average can acquire a nonzero average when multiplied by another term. Its contributions to
pressure and to the compatibility integrals in Equations~\eqref{eq:8.3} and \eqref{eq:8.15} must
therefore be retained at the next step.

\NSPage{100}%
For fixed base coefficients and stress, an actual velocity tuple
\NSi{NS-F-P100-001}{(\beta,v,\gamma,w)} determines the pressure and mean residuals in the
following order:
\begin{enumerate}[label=(\arabic*)]
\item Form \NSi{NS-F-P100-002}{W_{ab}=\langle w_aw_b\rangle_\theta}, using the full wave
velocities, and define \NSi{NS-F-P100-003}{g_r} by the third row of \eqref{eq:8.3}.
\item Define \NSi{NS-F-P100-004}{\mathcal P=\int\langle g_r\rangle_Y\,dR} and
\NSi{NS-F-P100-005}{p_m=\mathsf T_0(g_r-\rho\mathcal P)}.
\item With this pressure, define \NSi{NS-F-P100-006}{E_\theta,E_z} by the first two rows of
\eqref{eq:8.3}.
\item Define \NSi{NS-F-P100-007}{\mathcal J_\theta,\mathcal J_z} by \eqref{eq:8.15}, using the
same velocity and radial source.
\end{enumerate}
These definitions apply again after every change to the velocity. For either mean update, the new
tuple is
\NSFormula{NS-F-P100-008}{display}{none}
\[
(\beta,v,\gamma,w)_{\mathrm{new}}
=(\beta+\Delta\beta,v+\Delta v,\gamma+\Delta\gamma,w).
\]
For the temporal update, use the current \NSi{NS-F-P100-009}{E_\theta^\circ,E_z^\circ} in
\eqref{eq:8.20}; for the moment update, use the current
\NSi{NS-F-P100-010}{(\mathcal P,\mathcal J_\theta,\mathcal J_z)} in \eqref{eq:8.25}. In both
cases \NSi{NS-F-P100-011}{\Delta\gamma} is the actual axial increment from \eqref{eq:8.14}. In
particular, its cutoff remainder is included before the products in the four definitions above are
evaluated. After the moment update, Equations~\eqref{eq:8.26} to \eqref{eq:8.27} give the
recomputed defects exactly.

The correction cycle and its exponent estimates are given in Proposition~\ref{prop:9.6} and its
proof.

The definitions in physical variables ensure consistency across charts. The radial integrals keep
\NSi{NS-F-P100-012}{(z,t)} fixed, and \NSi{NS-F-P100-013}{q=q(z,t)} ensures that they introduce
no new dyadic bands. One common index can therefore be fixed on a neighborhood containing all
relevant support closures. The Fourier covariance \eqref{eq:8.23} proves agreement on overlaps; the
radial formulas agree by their physical definition. A change of representation uses only bounded
differences of covering indices. The integer-valued covering indices remain fixed during
differentiation, and the coordinate changes preserve the established derivative bounds. All estimates
above hold on the same domain \NSi{NS-F-P100-014}{0<q<q_*} fixed before the iteration. On a closed
region \NSi{NS-F-P100-015}{q\ge c>0}, including \NSi{NS-F-P100-016}{\eta=\pm1}, only finitely
many bands occur; the fixed-shift integrals, the Fourier inverse, and the finite-dimensional moment
correction preserve all one-sided source derivatives on this closed region, including the endpoint
derivatives.

\section{Residual improvement and the local field}\label{sec:9}
The wave and mean constructions in Sections~\ref{sec:7} and \ref{sec:8} supply corrections for the
different components of the momentum residual. We now combine them to prove Theorem~\ref{thm:3.1}.
Starting from the fixed background \NSi{NS-F-P100-017}{(u_B,p_B)} of
Proposition~\ref{prop:5.5} and the primary fields assembled in \eqref{eq:7.35} and made
divergence-free by Lemma~\ref{lem:7.7}, we must obtain a local velocity
\NSi{NS-F-P100-018}{u_{\mathrm{loc}}} whose residual and all its derivatives vanish to every order at
the singularity, while preserving the inner velocity growth and the exterior heat flow. Each
correction introduces new nonlinear errors, so we need a decay improvement after a complete cycle of
corrections.

A cycle first removes the nonzero angular harmonics by Proposition~\ref{prop:7.2}, then adjusts the
averaged stress by changing the primary amplitudes through \eqref{eq:7.31}. The mean equations
remove the part with zero auxiliary average through \eqref{eq:8.20}, and the five-equation system
\eqref{eq:8.25} corrects the three defects in compact support while preserving the two moment
constraints. We prove in Proposition~\ref{prop:9.6} that each cycle gains a fixed positive power of
\NSi{NS-F-P100-019}{\varepsilon}. All inverses retain the fixed background and primary amplitudes.
This permits every finite stage to be defined on one domain (Lemma~\ref{lem:9.7}), where the final
summation produces the local field of Proposition~\ref{prop:9.9}.

The notation guide in Table~\ref{tab:1} collects the recurring scales, operators, and field classes.
Throughout this section \NSi{NS-F-P100-020}{\kappa_s=10^{-5}}, and
\NSi{NS-F-P100-021}{\mathsf W_\alpha,\mathsf M_\alpha,\mathsf S_\alpha} have the meaning, at every derivative order,
given in Section~\ref{sec:6.4}. An exponent is always a power of
\NSi{NS-F-P100-022}{\varepsilon=Q^h} with the normalization
\NSPage{101}%
used for these classes. Constants, powers of \NSi{NS-F-P101-001}{S_*}, and inverse-distance powers
may depend on a fixed correction stage and a fixed order of amplitude differentiation, but are uniform
in band, label, copy, and point. The angular mean
\NSi{NS-F-P101-002}{\langle\cdot\rangle_\theta} and auxiliary mean
\NSi{NS-F-P101-003}{\langle\cdot\rangle_Y} remain distinct; their definitions are \eqref{eq:3.7}.
All class estimates refer to normalized chart coefficients. Physical velocity, pressure, and residual
are recovered with the factors \NSi{NS-F-P101-004}{Q^{-A}},
\NSi{NS-F-P101-005}{Q^{-2A}}, and \NSi{NS-F-P101-006}{Q^{-2A-1/2}}, respectively. In
particular, the residual of the physical fields is
\NSFormula{NS-F-P101-007}{display}{none}
\[
\mathscr R(u,p)=\partial_tu+(u\cdot\nabla)u-\Delta u+\nabla p,
\qquad
\mathscr R_*=Q^{2A+1/2}\mathscr R(u,p).
\]
We use physical fields in residual identities and normalized coefficients in the estimates below. The
Cartesian space--time jet notation \NSi{NS-F-P101-008}{|\cdot|_m} is defined in \eqref{eq:5.23}.

\subsection{Residual estimates for velocities constructed by curls.}\label{sec:9.1}
The amplitude equations \eqref{eq:7.5} cancel the principal part of a prescribed source. The
corresponding physical velocity also contains the terms introduced by taking a curl
(Lemma~\ref{lem:7.7}), and its residual contains derivatives of the amplitudes and background; we
will expand these terms in \eqref{eq:9.2}. We first estimate these remaining linear terms and the
nonlinear interactions in Proposition~\ref{prop:9.1} and Lemma~\ref{lem:9.2}, obtaining the gains
needed when successive corrections are added in Proposition~\ref{prop:9.6}.

We write an angular harmonic in a fixed label as
\NSi{NS-F-P101-009}{a\exp(ikm\Phi)}, where
\NSi{NS-F-P101-010}{m\in\mathbb Z\setminus\{0\}} and
\NSi{NS-F-P101-011}{a} is its amplitude. The phase normal
\NSi{NS-F-P101-012}{n_\Phi} is defined in \eqref{eq:7.4}, and the chain-rule operators for
evaluation at \NSi{NS-F-P101-013}{Y=Y(r,t)} are \eqref{eq:6.6}. The background coefficients
\NSi{NS-F-P101-014}{(b,V,G)} are those of Sections~\ref{sec:6} and \ref{sec:7}.

For a transverse amplitude \NSi{NS-F-P101-015}{t_m}, the potential in \eqref{eq:7.38} gives
\NSFormula{NS-F-P101-016}{display}{none}
\[
C_m=\frac{i n_\Phi\times t_m}{km|n_\Phi|^2},
\qquad
\operatorname{curl}_*(C_me^{ikm\Phi})=(t_m+r_m)e^{ikm\Phi},
\qquad
a_m=t_m+r_m.
\]
Here \NSi{NS-F-P101-017}{r_m} is the curl correction. If
\NSi{NS-F-P101-018}{t_m\in \mathsf W_\alpha}, then Lemma~\ref{lem:7.7} gives
\NSi{NS-F-P101-019}{r_m\in \mathsf W_{\alpha+1/2-\kappa_s}}. The longitudinal estimate \eqref{eq:7.39}
applies to the complete amplitude \NSi{NS-F-P101-020}{a_m}. This distinction matters in quadratic
transport: its advecting factor is the complete divergence-free velocity.

We separate the principal linear operator, which is handled by the amplitude equation, from the
remaining transport, pressure, and viscous terms. For later comparison with the residual, write the
principal amplitude-pressure operator as
\NSFormula{NS-F-P101-021}{display}{9.1}
\begin{equation}
\mathscr L_m(a,\pi)=Q^{1+h}N_{\mathrm{abs}}a+\mathcal Ka
+\varepsilon k^2m^2|n_\Phi|^2a+ikmn_\Phi\pi.
\tag{9.1}\label{eq:9.1}
\end{equation}
The inverse in Proposition~\ref{prop:7.2} constructs
\NSi{NS-F-P101-022}{(t_m,\pi_m)} satisfying
\NSi{NS-F-P101-023}{\mathscr L_m(t_m,\pi_m)=-f_m} and
\NSi{NS-F-P101-024}{n_\Phi\cdot t_m=0}, before the pulse cutoff is applied. The amplitude
\NSi{NS-F-P101-025}{t_m} is not required to have temporal zero extension before cutoff.

\begin{proposition}\label{prop:9.1}
Fix a label, \NSi{NS-F-P101-026}{m\in\mathbb Z\setminus\{0\}}, and
\NSi{NS-F-P101-027}{\alpha\in\mathbb R}. Let \NSi{NS-F-P101-028}{f_m\in \mathsf W_\alpha} satisfy the
support and smooth-extension hypotheses of Proposition~\ref{prop:7.2}. Assume also that the source
harmonic \NSi{NS-F-P101-029}{f_me^{ikm\Phi}} extends smoothly by zero outside the fixed local
rectangle, including its pulse endpoints. Let \NSi{NS-F-P101-030}{(t_m,\pi_m)} be the solution of
\NSi{NS-F-P101-031}{\mathscr L_m(t_m,\pi_m)=-f_m}. The same conclusions hold for a homogeneous
pulse \NSi{NS-F-P101-032}{t_m\in \mathsf W_\alpha}, with \NSi{NS-F-P101-033}{f_m=0}. Multiply its
potential and pressure by the pulse cutoff and take the curl as in \eqref{eq:7.38}. Then the
normalized linear residual, after adding the prescribed source
\NSi{NS-F-P101-034}{f_me^{ikm\Phi}}, belongs to
\NSi{NS-F-P101-035}{\mathsf W_{\alpha+1/2-3\kappa_s}}, apart from additive pulse-cutoff tails whose
Cartesian space--time derivatives vanish to infinite order as \NSi{NS-F-P101-036}{q\downarrow0}.
The pressure amplitude \NSi{NS-F-P101-037}{\pi_m} belongs to
\NSi{NS-F-P101-038}{\mathsf W_{\alpha+1/2}}.
\end{proposition}

\begin{proof}
We expand the linearized residual as its principal operator \eqref{eq:9.1} plus a remainder. For this
expansion, \NSi{NS-F-P101-039}{a} denotes an arbitrary harmonic amplitude and
\NSi{NS-F-P101-040}{\pi} its pressure amplitude. Let
\NSPage{102}%
\NSi{NS-F-P102-001}{E_{ik}=(\mathfrak t_*+bD_r+(V/R)\partial_\theta+GD_z)\Phi} be the phase
transport defect from \eqref{eq:7.9}. The remainder of the linearization about
\NSi{NS-F-P102-002}{(b,V,G)} is
\NSFormula{NS-F-P102-003}{display}{9.2}
\begin{equation}
\begin{aligned}
\mathscr L_m^{\mathrm{rem}}(a,\pi)={}&
 (-\varepsilon\partial_T+bD_r+GD_z)a+ikmE_{ik}a\\
&+a_r(\partial_Rb)e_r+a_zD_z(b,V,G)+(b/R)a_\theta e_\theta
 +(D_r\pi,0,D_z\pi)\\
&-\varepsilon\left\{
\begin{aligned}
 &(D_r^2+R^{-1}D_r+D_z^2)a-R^{-2}(a_r,a_\theta,0)\\
 &\quad+2ikm(n_{\Phi,r}D_r+n_{\Phi,z}D_z)a\\
 &\quad+ikm(D_rn_{\Phi,r}+n_{\Phi,r}/R+D_zn_{\Phi,z})a\\
 &\quad+2ikm(n_{\Phi,\theta}/R)(-a_\theta,a_r,0)
\end{aligned}
\right\}.
\end{aligned}
\tag{9.2}\label{eq:9.2}
\end{equation}
Indeed, transport of a cylindrical vector \NSi{NS-F-P102-004}{v} by
\NSi{NS-F-P102-005}{u} has the basis term
\NSi{NS-F-P102-006}{(-u_\theta v_\theta,u_\theta v_r,0)/R}. Linearizing it gives the connections
in \NSi{NS-F-P102-007}{\mathcal K} and the displayed \NSi{NS-F-P102-008}{b/R} term. The scalar
Laplacian of the harmonic contributes the quadratic phase-gradient term to \eqref{eq:9.1}, its two
mixed phase--amplitude derivative terms to \eqref{eq:9.2}, and the vector Laplacian contributes the
last angular connection and the \NSi{NS-F-P102-009}{R^{-2}} term. The amplitude coefficients are
independent of \NSi{NS-F-P102-010}{\theta}.

The transverse-amplitude equation of Proposition~\ref{prop:7.2} gives
\NSi{NS-F-P102-011}{\mathscr L_m(t_m,\pi_m)=-f_m}, with
\NSi{NS-F-P102-012}{\pi_m\in \mathsf W_{\alpha+1/2}}. The respective gains above
\NSi{NS-F-P102-013}{\alpha} in \eqref{eq:9.2} are
\NSFormula{NS-F-P102-014}{display}{none}
\[
\begin{array}{c|c}
\text{term}&\text{gain}\\ \hline
\text{slow transport}&1-\kappa_s\\
\text{phase transport defect}&1/2\\
\text{base derivatives and connections}&1\\
\text{pressure-amplitude gradient}&1/2-\kappa_s\\
\text{viscous amplitude derivatives}&1-2\kappa_s\\
\text{mixed derivatives and phase divergence}&1/2-\kappa_s\\
\text{viscous angular connection}&1/2
\end{array}
\]
The phase-defect bound permits polynomial factors of \NSi{NS-F-P102-015}{S_*}. Applying the
principal velocity operator to the curl correction preserves its exponent: fast time differentiates an
amplitude coefficient, and \NSi{NS-F-P102-016}{\varepsilon k^2=O(1)}. The potential formula and
these chain-rule operators give, at every derivative order, the common lower bound
\NSi{NS-F-P102-017}{\alpha+1/2-3\kappa_s} for the exponents of these terms.

The pulse cutoff is applied to the potential and pressure before the curl. The uncancelled source
\NSi{NS-F-P102-018}{(1-\psi)f}, and the principal pulse-cutoff derivative
\NSi{NS-F-P102-019}{\psi't} with its fixed multipliers, are supported where
\NSi{NS-F-P102-020}{P_v\le e^{-cS_*}}. For every fixed physical derivative their estimates have
the form \NSi{NS-F-P102-021}{CQ^{-M}S_*^Pe^{-cS_*}}, hence are flat because
\NSi{NS-F-P102-022}{S_*=\ell^2} and \NSi{NS-F-P102-023}{Q=2^{-\ell}}. Retain these terms in the
additive residual throughout the iteration; apply subsequent forward solves only to the supported
source terms.
\end{proof}

The nonlinear terms require separate estimates for interactions between waves and means and between
two waves. For two waves, incompressibility controls the contraction of the advecting amplitude with
the phase gradient.

For the next lemma, \NSi{NS-F-P102-024}{(w\cdot\nabla_*)v} denotes cylindrical transport with
normalized derivatives \NSi{NS-F-P102-025}{(D_r,R^{-1}\partial_\theta,D_z)}, including
differentiation of the cylindrical frame. Thus the lemma estimates the normalized nonlinear residual.

\begin{lemma}\label{lem:9.2}
Let \NSi{NS-F-P102-026}{w} be a wave in \NSi{NS-F-P102-027}{\mathsf W_\alpha}, and let
\NSi{NS-F-P102-028}{v} be a mean field with tangential components in
\NSi{NS-F-P102-029}{\mathsf M_\mu} and radial component in \NSi{NS-F-P102-030}{\mathsf M_{\mu+1}}. The nonzero
harmonics of \NSi{NS-F-P102-031}{(w\cdot\nabla_*)v+(v\cdot\nabla_*)w} belong to
\NSi{NS-F-P102-032}{\mathsf W_{\alpha+\mu-1/2}}.
